\section{An extremal horizontal defect and contact}\label{sec:envelope-contact}

We next turn a nonaffine member of $\Aclass_{m,C_0}$ into a
contact between a solution and a subsolution. A finite horizontal
probability configuration consists of positions
$y_1,\ldots,y_N\in\R^m$ and weights
$\theta_i>0$ with $\sum_i\theta_i=1$. Its mean and root mean
square radius are
\[
  x=\sum_i\theta_i y_i,\qquad
  r=\left(\sum_i\theta_i|y_i-x|^2\right)^{1/2}.
\]

\subsection{The barycentric envelope}

Let $I\subseteq\R$ be an open interval. For a continuous function $f$ on
$\R^m\times I$ with subquadratic spatial growth on compact time intervals,
define, for $s>0$,
\begin{equation}\label{eq:barycentric-envelope}
  \Cenv_s f(x,t)=
  \sup_{\substack{N,\ \theta_i>0,\ \sum_i\theta_i=1\\
                  \sum_i\theta_i y_i=x}}
   \left\{\sum_i\theta_i f(y_i,t)
             -\frac{\sum_i\theta_i|y_i-x|^2}{2s}\right\}.
\end{equation}
The supremum is over finite configurations, including the
single point $y_1=x$.
This operation averages one horizontal slice at the fixed time $t$;
$s$ controls the variance penalty. As the next lemma shows,
$\Cenv_s f(\cdot,t)-|\cdot|^2/(2s)$ is the least concave majorant of
$f(\cdot,t)-|\cdot|^2/(2s)$.

For the following estimate, suppose additionally that some $\kappa\ge0$
satisfies
\begin{equation}\label{eq:jensen-bound}
  \sum_i\theta_i f(y_i,t)-f(x,t)\le\kappa r
\end{equation}
for every $t\in I$ and every finite horizontal probability
configuration, with mean $x$ and radius $r$. Since
\[
  \kappa r-\frac{r^2}{2s}
   =\frac{\kappa^2s}{2}-\frac{(r-\kappa s)^2}{2s},
\]
the definition gives
\begin{equation}\label{eq:barycentric-sandwich}
  0\le\Cenv_s f(x,t)-f(x,t)
     \le\sup_{r\ge0}\left(\kappa r-\frac{r^2}{2s}\right)
     =\frac{\kappa^2s}{2}.
\end{equation}
The scalar maximum occurs at $r=\kappa s$. The extremal construction
below supplies a solution $v$ and a positive $\kappa$ attaining this bound.
We will prove that $\Cenv_t v-\kappa^2t/2$ is a subsolution for $t>0$;
the sandwich places it below $v$, with contact where equality holds and
a strictly positive unshifted gap.
The additional Jensen assumption applies only to this estimate; the next
lemma uses only the stated continuity and growth hypotheses.

\begin{lemma}[Barycentric envelope]\label{lem:barycentric-envelope}
Let $f$ have the preceding growth property on $\R^m\times I$.
For $s>0$ the supremum in \eqref{eq:barycentric-envelope} is
attained by a finite configuration, and
\begin{equation}\label{eq:barycentric-duality}
  \Cenv_s f(x,t)-\frac{|x|^2}{2s}
   =\min_{A\in\R^m}\ \max_{y\in\R^m}
       \left\{f(y,t)-\frac{|y|^2}{2s}-A\cdot(y-x)\right\}.
\end{equation}
Both extrema on the right are attained. Allowing any probability
measure of finite second moment and mean $x$ in
\eqref{eq:barycentric-envelope} gives the same value.
The function $\Cenv_s f(x,t)$ is continuous in $(x,t,s)$ for
$s>0$, and
\begin{equation}\label{eq:barycentric-semiconcavity}
  x\longmapsto\Cenv_s f(x,t)-\frac{|x|^2}{2s}
       \quad\text{is concave}.
\end{equation}
If $f_j\in C(\R^m\times I)$, $f_j\to f$ locally uniformly, and the sequence has the common
growth control
\[
  \lim_{M\to\infty}
  \sup_j\sup_{\substack{|y|\ge M\\t\in J}}
       \frac{|f_j(y,t)|}{|y|^2}=0
       \qquad(J\Subset I),
\]
then $\Cenv_s f_j\to\Cenv_s f$ locally uniformly for $s>0$.
\end{lemma}

\begin{proof}
For fixed $(t,s)$ put $g(y)=f(y,t)-|y|^2/(2s)$ and
\[
  \Gamma(A;x,t,s)=\sup_y\{g(y)-A\cdot(y-x)\}.
\]
The expression on the left of \eqref{eq:barycentric-duality}
is the supremum of $\sum_i\theta_i g(y_i)$ over configurations
of mean $x$, because
$\sum_i\theta_i|y_i-x|^2=\sum_i\theta_i|y_i|^2-|x|^2$.
It is therefore at most $\inf_A\Gamma(A;x,t,s)$.

Fix a compact range of $(x,t,s)$ with $s>0$. For $A\ne0$, use
$y=x-A/|A|$ in the supremum to get
$\Gamma(A;x,t,s)\ge |A|-C$ on this range. The value at $A=0$
is uniformly bounded above, since the negative quadratic part
of $g$ dominates the subquadratic growth of $f$. Thus every
minimizing sequence whose values are at most $\Gamma(0)+1$
lies in a fixed bounded set of slopes.
For slopes in that set, the quantity
$g(y)-A\cdot(y-x)$ tends uniformly to $-\infty$ as
$|y|\to\infty$. Its maximizing positions therefore lie in one
compact set. It follows that $\Gamma$ is continuous on the
relevant set, the minimum and maxima are attained, and their
value depends continuously on $(x,t,s)$.

Let $A_0$ be a minimizing slope, and let $E$ be its nonempty
compact set of maximizing positions. Then $x$ belongs to the
convex hull of $E$. Otherwise strict separation gives a vector
$e$ and $\delta>0$ with
$e\cdot(y-x)\ge\delta$ for all $y\in E$. If $y_h$ maximizes
$\Gamma(A_0+he;x,t,s)$ for $h>0$, the uniform compactness of
maximizers and continuity show that every limit of $y_h$ as
$h\downarrow0$ belongs to $E$. For small $h$,
$e\cdot(y_h-x)\ge\delta/2$, and hence
\[
  \Gamma(A_0+he;x,t,s)
   \le\Gamma(A_0;x,t,s)-h\delta/2,
\]
contradicting minimality. The finite-dimensional convex hull
theorem now represents $x$ as a convex combination of finitely
many points of $E$. Discard any zero weights. Their weighted
values of $g$ equal $\Gamma(A_0;x,t,s)$, proving the reverse
inequality in \eqref{eq:barycentric-duality} and finite
attainment in \eqref{eq:barycentric-envelope}.

The affine majorant
$g(y)\le\Gamma(A_0;x,t,s)+A_0\cdot(y-x)$ can also be
integrated against any probability measure of mean $x$ and
finite second moment. All terms are integrable by the growth
hypothesis. Thus such measures cannot increase the value.
The right side of \eqref{eq:barycentric-duality} is an infimum
of affine functions of $x$, so it is concave in $x$. This is
\eqref{eq:barycentric-semiconcavity}.

For the convergence assertion, on each compact range of
$(x,t,s)$ the same preceding bounds on slopes and maximizing
positions hold uniformly for $f_j$: local uniform convergence
gives common bounds near the compact range, and the common
subquadratic control handles large $y$. Both extrema in
\eqref{eq:barycentric-duality} can therefore be restricted to
fixed compact sets. On these sets the expressions converge
uniformly, proving the asserted local uniform convergence.
\end{proof}

\subsection{A common-time subsolution rule}

To test the barycentric envelope, we must combine subsolution inequalities
at several spatial points sharing one time variable. Lemma~\ref{lem:sum-rule} is
applied after separating and then synchronizing those times; the penalty
Hessian vanishes on the direction with time speed $1$ in every component,
which is the direction used by $F$.

\begin{lemma}[Common-time subsolutions]\label{lem:common-time-subsolution}
Let $v$ be a continuous viscosity subsolution of
\eqref{eq:limit-pde} on $\R^m\times I$, and fix positive weights
$\theta_i$ with $\sum_{i=1}^N\theta_i=1$. Define
\[
  V((y_i)_{i=1}^N,t)=\sum_{i=1}^N\theta_i v(y_i,t).
\]
If a $C^2$ function $\Phi$ touches $V$ from above at $(Y_0,t_0)$,
where $Y=(y_i)_{i=1}^N$, then
\begin{equation}\label{eq:common-time-inequality}
  D^2\Phi(Y_0,t_0)
   \left[\left(
       \bigl(\theta_i^{-1}\nabla_{y_i}\Phi(Y_0,t_0)\bigr)_{i=1}^N,
       1\right)\right]\ge0.
\end{equation}
\end{lemma}

\begin{proof}
Replace $\Phi$ by
$\Phi^\sharp=\Phi+|Y-Y_0|^4+|t-t_0|^4$. This makes the local
maximum strict while preserving the derivatives through order
two at the contact. Choose a small closed product neighborhood
of $((y_i^0,t_0))_{i=1}^N$, contained in the domain, such that
the diagonal points $(Y,t)$ lie in this strict neighborhood.
For independent time variables $T=(t_1,\ldots,t_N)$, put
\[
  \bar t=\frac1N\sum_i t_i,\qquad
  \chi(T)=\frac12\sum_i(t_i-\bar t)^2,\qquad
  \Psi_\lambda(Y,T)=\Phi^\sharp(Y,\bar t)+\lambda\chi(T).
\]
Maximize
$\sum_i\theta_i v(y_i,t_i)-\Psi_\lambda(Y,T)$ on the closed
product. Comparison with the original diagonal point and
boundedness on this product show
$\lambda\chi(T_\lambda)\le C$ at any maximum. Thus all
time deviations tend to zero as $\lambda\to\infty$. Every
limit of such maximum points is diagonal; dropping the
nonnegative penalty in the maximum inequality shows that its
diagonal limit is a maximum of $V-\Phi^\sharp$. Strictness
forces this limit to be $(Y_0,t_0)$. Hence all the maximum
points converge to the original point and are interior for
sufficiently large $\lambda$.

Fix one such $\lambda$ and one maximum. Set
\[
  q_i=\theta_i^{-1}\nabla_{y_i}
             \Phi^\sharp(Y_\lambda,\bar t_\lambda),
  \qquad
  \zeta_\lambda=((q_i,1))_{i=1}^N,\qquad
  A_\lambda=D^2\Psi_\lambda(Y_\lambda,T_\lambda),
\]
where $A_\lambda$ is written in the separate variables
$((y_i,t_i))_i$. Apply Lemma~\ref{lem:sum-rule} to the
functions $\theta_i v(y_i,t_i)$ and the test $\Psi_\lambda$.
For every $\eta>0$ it supplies closed superjets with Hessians
$X_i$ and first derivatives $D_{(y_i,t_i)}\Psi_\lambda$, with
$\operatorname{diag}(X_i)\le A_\lambda+\eta A_\lambda^2$.
Dividing each jet by the positive number $\theta_i$ gives a
closed superjet of $v$ with horizontal gradient $q_i$.
The viscosity inequalities therefore give
$X_i[(q_i,1)]\ge0$. Evaluating the matrix inequality on
$\zeta_\lambda$ yields
\[
  0\le\sum_iX_i[(q_i,1)]
     \le A_\lambda[\zeta_\lambda]
                   +\eta|A_\lambda\zeta_\lambda|^2.
\]
First let $\eta\downarrow0$, keeping $\lambda$ and its maximum
fixed. The quantities $A_\lambda$ and $\zeta_\lambda$ do not
depend on $\eta$, so $A_\lambda[\zeta_\lambda]\ge0$.
The time component of $\zeta_\lambda$ is $(1,\ldots,1)$,
which is in the kernel of $D^2\chi$. The linear map
$(Y,T)\mapsto(Y,\bar t)$ sends $\zeta_\lambda$ to
$((q_i)_i,1)$. We conclude that
\[
  D^2\Phi^\sharp(Y_\lambda,\bar t_\lambda)
       [((q_i)_i,1)]\ge0.
\]
Now let $\lambda\to\infty$. The maximum points converge to the
contact, and the horizontal derivatives converge there.
The derivatives of $\Phi^\sharp$ at that point are those of
$\Phi$, giving \eqref{eq:common-time-inequality}.
The time components of the first derivatives of
$\Psi_\lambda$ need no bound: they do not enter the direction
$(\nabla_xv,1)$ in the individual viscosity inequalities.
\end{proof}

\begin{corollary}[A barycentric subsolution]\label{cor:barycentric-subsolution}
Let $f\in C(\R^m\times(0,\infty))$ be a viscosity subsolution of
\eqref{eq:limit-pde} with subquadratic spatial growth on compact time
intervals. Then $(x,t)\mapsto\Cenv_t f(x,t)$ is a viscosity subsolution of
\eqref{eq:limit-pde} on $\R^m\times(0,\infty)$. For every $c\in\R$, the same
holds for $(x,t)\mapsto\Cenv_t f(x,t)-ct$.
\end{corollary}

\begin{proof}
Continuity follows from Lemma~\ref{lem:barycentric-envelope}. Fix $c\in\R$,
and let $\phi$ touch $(x,t)\mapsto\Cenv_t f(x,t)-ct$ from above at
$(x_0,t_0)$ with $t_0>0$. Choose an attaining finite configuration
$(\theta_i,y_i^0)$ for $\Cenv_{t_0}f(x_0,t_0)$, and keep its positive weights
fixed. For variable $Y=(y_i)_i$ let $X(Y)=\sum_i\theta_i y_i$ and define
\[
  \Phi(Y,t)=\phi(X(Y),t)+ct
              +\frac{\sum_i\theta_i|y_i-X(Y)|^2}{2t}.
\]
The defining supremum for $\Cenv_t$ shows that
$\sum_i\theta_i f(y_i,t)-\Phi(Y,t)$ has a local maximum at $(Y_0,t_0)$:
for nearby $(Y,t)$ it is at most
$\Cenv_t f(X(Y),t)-ct-\phi(X(Y),t)$, and equality holds at the chosen
configuration. Set $p=\nabla_x\phi(x_0,t_0)$. The individual horizontal
directions in Lemma~\ref{lem:common-time-subsolution} are
\[
  q_i=\theta_i^{-1}\nabla_{y_i}\Phi
       =p+\frac{y_i^0-x_0}{t_0}.
\]
Their weighted mean is $p$. Along the straight motion with velocities
$y_i'=q_i$ and $t'=1$, the mean has velocity $p$ and each offset $y_i-X$
has velocity $(y_i^0-x_0)/t_0$. Consequently, along this motion the variance
term is
\[
  \frac{\sum_i\theta_i
           |(y_i^0-x_0)(1+h/t_0)|^2}{2(t_0+h)}
    =\frac{t_0+h}{2t_0^2}
           \sum_i\theta_i|y_i^0-x_0|^2,
\]
which is affine in $h$. The term $ct$ is also affine, and the Hessian of the
remaining term along this motion is $D^2\phi[(p,1)]$.
Lemma~\ref{lem:common-time-subsolution}, applied to the subsolution
inequality for $f$, therefore gives $F[\phi](x_0,t_0)\ge0$. This proves the
claim for every $c$, including $c=0$.
\end{proof}

\subsection{Construction and slope of the contact}

\begin{proposition}[Extremal envelope contact]\label{prop:envelope-contact}
If $\Aclass_{m,C_0}$ contains a nonaffine member, then there
exist $\kappa\in(0,\infty)$ and $v\in\Aclass_{m,C_0}$ centered
at $z_*=(0,1)$ for which \eqref{eq:jensen-bound} holds with $f=v$ and
$I=\R$.
For $t>0$ set
\begin{equation}\label{eq:contact-envelope}
  W(x,t)=\Cenv_t v(x,t)-\frac{\kappa^2t}{2}.
\end{equation}
This function is continuous, has subquadratic spatial growth
on compact positive-time intervals, and is a viscosity
subsolution of \eqref{eq:limit-pde} on $\R^m\times(0,\infty)$.
It satisfies
\begin{equation}\label{eq:contact-sandwich}
  v(x,t)-\frac{\kappa^2t}{2}\le W(x,t)\le v(x,t)
       \quad(t>0),\qquad W(z_*)=v(z_*)=0,
\end{equation}
and $W(\cdot,t)-|x|^2/(2t)$ is concave for every $t>0$.
At every point of contact with $t>0$, the barycentric gap satisfies
\begin{equation}\label{eq:contact-barycentric-gap}
  W(x,t)=v(x,t)\quad\Longrightarrow\quad
  \frac{\Cenv_t v(x,t)-v(x,t)}{t}=\frac{\kappa^2}{2}>0.
\end{equation}
Finally,
\begin{equation}\label{eq:contact-slope}
  H_v(z_*)\in\R,\qquad
  \lim_{\ell\downarrow0}\Senv_\ell^+W(z_*)=H_v(z_*).
\end{equation}
\end{proposition}

\begin{proof}
\emph{A positive finite horizontal defect.}
Start with a nonaffine member $v^{(0)}$ and center it at the
origin using Lemma~\ref{lem:frame-invariance}. Some horizontal
slice of $v^{(0)}$ is nonaffine. To verify this, suppose instead
that
$v^{(0)}(x,t)=p(t)\cdot x+c(t)$ for every $t$. The coefficients
$p,c$ are continuous, as follows by evaluating the function at
$0$ and at the coordinate vectors. Direct optimization gives
\[
  \Penv_\ell v^{(0)}(x,t)
   =p(t+\ell)\cdot x+c(t+\ell)
                         +\frac{\ell}{2}|p(t+\ell)|^2.
\]
For every fixed $x,t$ this is convex in $\ell$. In each
convexity inequality, the coefficient of $x$ must vanish,
since the inequality holds for all $x\in\R^m$. Hence $p$
satisfies the affine interpolation identity on every interval,
and so is affine in $t$. On the other hand,
\eqref{eq:centered-growth}, applied at $x=\pm Ke$ and
$|t|\le K$, implies $|p(t)|\le C L(K)^2$. An affine function
of $t$ with this bound as $K\to\infty$ is constant. For
constant $p$, the displayed forward formula and its backward
counterpart
\[
  \Menv_\ell v^{(0)}(x,t)
   =p\cdot x+c(t-\ell)-\frac{\ell}{2}|p|^2
\]
show that $c$ is both convex and concave. It is therefore
affine, contradicting the choice of $v^{(0)}$.

A continuous nonaffine spatial slice has a nonzero midpoint
Jensen defect, since the midpoint identity everywhere would
make that slice affine. The transformation
$f(x,t)\mapsto-f(x,-t)$ preserves $\Aclass_{m,C_0}$: it
exchanges the two viscosity inequalities, and it sends a fit
$(p,a)$ to $(-p,a)$. It also changes the sign of a horizontal
Jensen defect. Apply this transformation if needed, so that
$v^{(0)}$ has a positive defect. Define
\begin{equation}\label{eq:jensen-supremum}
  \kappa=\sup_{\substack{t,\ \text{finite configurations}\\r>0}}
     \frac{\sum_i\theta_i v^{(0)}(y_i,t)-v^{(0)}(x,t)}{r}.
\end{equation}
Then $\kappa>0$. It is finite as well. For any configuration in
this supremum, choose a fit at its mean $(x,t)$ of radius $r$
and use \eqref{eq:frame-change} with these data. At time zero
the new positions are $(y_i-x)/r$, with mean zero and second
moment one. The defect divided by $r$ is unchanged, since the
affine terms cancel in the weighted difference. The new
function is centered. Its growth bound implies
$|\bar v(y,0)|\le C(1+|y|^2)$, with one constant for all these
normalizations. Integrating this inequality over a
configuration of second moment one bounds the quotient by a
fixed constant. Thus $0<\kappa<\infty$. More generally,
\eqref{eq:frame-change} preserves every such quotient, at
every time: spatial differences and their root mean square
radius both scale by $R$, and the subtracted affine terms
cancel.

\emph{Attainment after compactness.}
Choose configurations whose quotients in
\eqref{eq:jensen-supremum} tend to $\kappa$, and for each one
make the normalization just described. Denote the resulting
functions by $v_j$ and the normalized finite probability
measures by $\mu_j$. Then
\[
  v_j(0,0)=0,\qquad
  \int y\,d\mu_j=0,\qquad
  \int |y|^2\,d\mu_j=1,\qquad
  \int v_j(y,0)\,d\mu_j\longrightarrow\kappa.
\]
Every $v_j$ is centered at the origin, and every finite
configuration for it satisfies the Jensen bound with constant
$\kappa$. First use Lemma~\ref{lem:fit-compactness} to pass to
a locally uniform limit $\widehat v\in\Aclass_{m,C_0}$,
centered at the origin. The Jensen bound passes to this limit
for each fixed finite configuration.

The moment bound gives
$\mu_j(\{|y|>M\})\le M^{-2}$, so the measures are tight.
After a further subsequence they converge weakly to a
probability measure $\mu$. Lower semicontinuity of the second
moment gives $\int|y|^2\,d\mu\le1$. The mean remains zero:
the bound
$\int_{|y|>M}|y|\,d\mu_j\le M^{-1}$ makes the first moments
uniformly integrable, and the same bound holds for the limit.
To retain the extremal defect, we also need convergence of the value
integrals despite the varying functions and measures. For a compactly
supported continuous cutoff $\chi_M$, local uniform
convergence and weak convergence give
\[
  \int\chi_M(y)v_j(y,0)\,d\mu_j
       \longrightarrow
  \int\chi_M(y)\widehat v(y,0)\,d\mu.
\]
Choose $\chi_M=1$ on $\{|y|\le M\}$ and zero outside
$\{|y|<2M\}$. By
\eqref{eq:uniform-subquadratic-growth},
\[
  \int_{|y|>M}|v_j(y,0)|\,d\mu_j
  \le
  \sup_{\substack{j\\|y|>M}}\frac{|v_j(y,0)|}{|y|^2}
       \int|y|^2\,d\mu_j
  \longrightarrow0
\]
uniformly in $j$; the analogous tail bound holds for
$\widehat v$ and $\mu$. Letting first $j\to\infty$ and then
$M\to\infty$ proves
\begin{equation}\label{eq:extremal-measure}
  \int y\,d\mu=0,\qquad
  \int|y|^2\,d\mu\le1,\qquad
  \int\widehat v(y,0)\,d\mu-\widehat v(0,0)=\kappa.
\end{equation}

Apply \eqref{eq:barycentric-sandwich} to $\widehat v$. By
Lemma~\ref{lem:barycentric-envelope}, the measure in
\eqref{eq:extremal-measure} is admissible for the value of
$\Cenv_{1/\kappa}\widehat v(0,0)$. It gives
\[
  \Cenv_{1/\kappa}\widehat v(0,0)-\widehat v(0,0)
     \ge\kappa-\frac{\kappa}{2}\int|y|^2\,d\mu
     \ge\frac{\kappa}{2}.
\]
This is equality by \eqref{eq:barycentric-sandwich}. In
particular, loss of second moment in the weak limit causes no
loss of contact; the equality also forces
$\int|y|^2\,d\mu=1$.

Set $R=1/\kappa$, and choose a fit $(p_0,a_0)$ for
$\widehat v$ at $(0,0)$ of radius $R$. Define the final function
by the single change of frame and time translation
\[
  v(y,t)=
   \frac{\widehat v(R(y+p_0(t-1)),R(t-1))-\widehat v(0,0)}{R}
       -p_0\cdot y
       -\left(a_0+\frac{|p_0|^2}{2}\right)(t-1).
\]
It is centered at $(0,1)$ and still has the global Jensen bound
with constant $\kappa$. Under the change of frame
\eqref{eq:frame-change}, at corresponding points $Z$ and
$(y,s)$ the envelope gaps obey
\[
  \Cenv_\sigma\bar v(y,s)-\bar v(y,s)
     =\frac{\Cenv_{R\sigma}\widehat v(Z)-\widehat v(Z)}{R}.
\]
This follows by scaling the spatial variance by $R^2$ in
\eqref{eq:barycentric-envelope}; the affine terms again
cancel. The preceding equality at parameter $1/\kappa=R$
therefore becomes
\begin{equation}\label{eq:normalized-envelope-contact}
  \Cenv_1v(0,1)-v(0,1)=\frac{\kappa^2}{2}.
\end{equation}

\emph{The subsolution below $v$.}
Define $W$ by \eqref{eq:contact-envelope}.
Equations \eqref{eq:barycentric-sandwich} and
\eqref{eq:normalized-envelope-contact} give
\eqref{eq:contact-sandwich}. At any contact, \eqref{eq:contact-envelope}
also gives \eqref{eq:contact-barycentric-gap}.
Lemma~\ref{lem:barycentric-envelope}
gives continuity and the stated spatial semiconcavity. The
sandwich and \eqref{eq:uniform-subquadratic-growth} give
subquadratic spatial growth on every compact positive-time
interval.
Since $v$ is a solution with the growth in
\eqref{eq:uniform-subquadratic-growth},
Corollary~\ref{cor:barycentric-subsolution}, with $f=v$ and
$c=\kappa^2/2$, gives the subsolution property of $W$.

\emph{A finite slope at the contact.}
Lemma~\ref{lem:envelope-slopes} already bounds $H_v$ above at every point.
We first bound it below at the finitely many points supporting the extremal
configuration, using one common forward paraboloid, and then transfer this
bound to their barycenter through equality in the Jensen bound.
Choose a finite attaining configuration in
\eqref{eq:normalized-envelope-contact}, with mean $x=0$,
radius $r$, and positive weights. If
$D=\sum_i\theta_i v(y_i,1)-v(0,1)$, its attainment and the
Jensen bound imply
\[
  \frac{\kappa^2}{2}=D-\frac{r^2}{2}
       \le\kappa r-\frac{r^2}{2}
       =\frac{\kappa^2}{2}-\frac{(r-\kappa)^2}{2}.
\]
All inequalities are equalities, so
\begin{equation}\label{eq:contact-configuration}
  r=\kappa,\qquad
  \sum_i\theta_i v(y_i,1)-v(0,1)=\kappa r.
\end{equation}
Let $A$ be a minimizing slope in
\eqref{eq:barycentric-duality} for $(x,t,s)=(0,1,1)$.
The weighted attaining value equals its affine majorant.
Since all weights are positive, each $y_i$ must attain that
majorant. The identity
\[
  v(y,1)-\frac{|y|^2}{2}-A\cdot y
   =v(y,1)-\frac{|y+A|^2}{2}+\frac{|A|^2}{2}
\]
therefore shows that all $y_i$ are optimizers for $\Penv_1v(b,0)$ with
$b=-A$.

Put $h(s)=\Penv_s v(b,0)$ for $0\le s\le1$. It is finite and
convex by Lemma~\ref{lem:quadratic-envelopes}. For
$0<\ell<1$, use the point
$y_{i,\ell}^-=b+(1-\ell)(y_i-b)$ at time $1-\ell$ as a competitor for
$\Menv_\ell v(y_i,1)$. Its value is bounded above using
$\Penv_{1-\ell}v(b,0)$. The two quadratic costs on this linear
path add, giving explicitly
\[
  \Menv_\ell v(y_i,1)
      \le h(1-\ell)+\frac{|y_i-b|^2}{2},
  \qquad
  v(y_i,1)=h(1)+\frac{|y_i-b|^2}{2}.
\]
It follows that
\[
  \Senv_\ell^-v(y_i,1)
     \ge\frac{h(1)-h(1-\ell)}{\ell}
     \ge h(1)-h(0).
\]
The last inequality is the chord inequality for the convex
function $h$. Taking the slope limit shows
$H_v(y_i,1)\ge h(1)-h(0)>-\infty$. These finitely many values
are also finite above by Lemma~\ref{lem:envelope-slopes}.

For each $\ell>0$, choose an optimizer $\widehat y_i$ for
$\Penv_\ell v(y_i,1)$ and let
$\widehat x=\sum_i\theta_i\widehat y_i$. Write
$\xi_i=y_i-x$, $\widehat\xi_i=\widehat y_i-\widehat x$, and
$r'=(\sum_i\theta_i|\widehat\xi_i|^2)^{1/2}$. The exact
orthogonal decomposition and the reverse triangle inequality
in the weighted Euclidean space give
\[
  \sum_i\theta_i|\widehat y_i-y_i|^2
    =|\widehat x-x|^2+
        \sum_i\theta_i|\widehat\xi_i-\xi_i|^2
    \ge|\widehat x-x|^2+(r'-r)^2.
\]
Use the Jensen bound at time $1+\ell$,
\eqref{eq:contact-configuration} at time $1$, and
$\widehat x$ as a forward competitor from $x$. The
decomposition yields
\begin{align*}
  \ell\sum_i\theta_i\Senv_\ell^+v(y_i,1)
   &\le \ell\Senv_\ell^+v(x,1)
       +\kappa(r'-r)-\frac{(r'-r)^2}{2\ell}\\
   &\le \ell\Senv_\ell^+v(x,1)+\frac{\kappa^2\ell}{2}.
\end{align*}
Divide by $\ell$ and take $\ell\downarrow0$. There are only
finitely many endpoints, all with finite limits, so
\[
  H_v(z_*)\ge
     \sum_i\theta_i H_v(y_i,1)-\frac{\kappa^2}{2}
     \ge h(1)-h(0)-\frac{\kappa^2}{2}>-\infty.
\]
Its upper finiteness is part of
Lemma~\ref{lem:envelope-slopes}; hence $H_v(z_*)\in\R$.

It remains to identify the forward slope of $W$. At the
contact, the sandwich gives for every $\ell>0$
$\Senv_\ell^+W(z_*)\le\Senv_\ell^+v(z_*)$. For
$0<\ell<1$ it also gives
$\Senv_\ell^-W(z_*)\ge\Senv_\ell^-v(z_*)$, since the
backward time is positive. The subsolution property and growth
of $W$ allow the one-sided inequality
\eqref{eq:subsolution-slope-order} to be applied to it.
Consequently
\[
  H_v(z_*)\le
  \limsup_{\ell\downarrow0}\Senv_\ell^-W(z_*)
  \le\lim_{\ell\downarrow0}\Senv_\ell^+W(z_*)
  \le H_v(z_*).
\]
This proves \eqref{eq:contact-slope}.
\end{proof}
